\documentclass[article,shortnames,nojss]{jss}

%% --- packages ---
\usepackage{amsmath,amssymb,amsthm}
\usepackage{booktabs}
\usepackage{array}
\usepackage{graphicx}
\usepackage{lmodern}
\usepackage{microtype}
\usepackage{orcidlink}

%% --- overflow tolerance (long code identifiers break the right margin) ---
\setlength{\emergencystretch}{4em}
\hbadness=10000
\hfuzz=20pt

%% --- math macros ---
\providecommand{\Rset}{\mathbb{R}}
\renewcommand{\E}{\mathbb{E}}
\newcommand{\indep}{\perp\!\!\!\perp}
\newcommand{\OTIS}{\textsf{OTIS}}
\newcommand{\SIU}{\textsf{SIU}}
\newcommand{\TPS}{\textsf{TPS}}

\title{morie: Multi-domain Open Research and Inferential Estimation in \proglang{Python}}
\Plaintitle{morie: Multi-domain Open Research and Inferential Estimation in Python}
\Shorttitle{morie: Multi-domain Inferential Estimation in Python}

\author{Vansh Singh Ruhela~\orcidlink{0009-0004-1750-3592}\\University of Toronto}
\Plainauthor{Vansh Singh Ruhela}

\Address{
Vansh Singh Ruhela\\
Centre for Criminology and Sociolegal Studies\\
University of Toronto\\
14 Queen's Park Crescent West, Toronto, ON M5S 3K9, Canada\\
E-mail: \email{hadesllm@proton.me}\\
ORCID: \orcidlink{0009-0004-1750-3592}\\
Pronunciation: vansh sing roo-HAY-lah
}

\Abstract{
The \proglang{Python} package \pkg{morie} (Multi-domain Open
Research and Inferential Estimation) implements a curated
dual-language toolkit for observational inference, with explicit
support for the MRM (Multilevel Reconciliation Methodology) framework as a
sociolegal application target. The package provides a unified
entry point to double-machine-learning interactive regression and
partial-linear estimators by wrapping the \proglang{Python}
package \pkg{DoubleML} \citep{Bach2022DoubleMLpy} via the
\pkg{scikit-learn} learner interface \citep{Pedregosa2011}; it
provides independent implementations of inverse-probability-weighting,
augmented inverse-probability-weighting, propensity-score
matching, survey-weighted estimation, signal processing
primitives, spatial statistics, statistical physics of crime
(Hawkes processes), psychometric estimators, and a suite of MRM
modules for ingestion of Canadian open carceral data. The package
is released on the Python Package Index (PyPI) at
\href{https://pypi.org/project/morie/}{pypi.org/project/morie} with
PEP 740 sigstore attestations on every wheel and source
distribution. This paper describes the \proglang{Python}
implementation, the wrapper API around \pkg{DoubleML}, the
\pkg{RichResult} return-type design, and three reproducible
examples on the open Ontario, federal, and Toronto Police Service
records that the MRM framework integrates.  Version~0.5.0 ships
275 textbook-derived callables in \code{morie.fn} across
fourteen suites with full \proglang{R} parity, and adopts a
per-component licensing model (\proglang{Python}:
\code{MIT OR Apache-2.0}; \proglang{R}: \code{GPL-2.0-only}).
}
\Keywords{causal inference, observational inference,
double machine learning, sociolegal statistics, Hawkes process,
psychometrics, \proglang{Python}}
\Plainkeywords{causal inference, observational inference,
double machine learning, sociolegal statistics, Hawkes process,
psychometrics, Python}
% --- morie version stamp -----------------------------------------------------
% This paper documents \pkg{morie} v0.5.0 (released 2026-05-13).  Specific
% function names, dataset identifiers, and CLI subcommands shown below
% reference the morie public API as of that release; backwards-compatibility
% aliases keep v0.x code references resolving across the v0.4.x series.
% -----------------------------------------------------------------------------


\begin{document}
\sloppy

\section{Introduction}
\label{sec:intro}

The \proglang{Python} package \pkg{morie} provides a curated
multi-domain toolkit for observational inference. The Python
implementation is dual-licensed under \code{MIT OR Apache-2.0}
(the Rust-ecosystem convention; recipient picks either) so that
\pkg{morie} can be used freely in any downstream Python project
regardless of that project's own licensing.  The R companion
package \citep{Ruhela2026MorieR} is released under
GPL-2.0-only to match the R-ecosystem and CRAN convention.  The
package supports a
unified estimator interface for the average treatment effect
(ATE), the average treatment effect on the treated (ATT), the
conditional and group average treatment effects (CATE, GATE), the
local average treatment effect (LATE), and the augmented
inverse-probability-weighting (AIPW) estimator. Double machine
learning is exposed through wrappers around the \pkg{DoubleML}
classes \citep{Bach2022DoubleMLpy, Bach2024DoubleML}, which in
turn rely on \pkg{scikit-learn} \citep{Pedregosa2011} for the
nuisance-function learners. Outside causal inference, the package
provides signal processing primitives, cryptographic primitives,
spatial statistics, statistical-physics models of crime, and
psychometric estimators that complement the causal-inference core.

\pkg{morie} is positioned alongside the established
\proglang{Python} ecosystem of causal-inference packages:
\pkg{DoubleML} \citep{Bach2022DoubleMLpy} for canonical
double-machine-learning estimators, \pkg{causalml} (Uber) for
heterogeneous treatment effects, \pkg{EconML} (Microsoft) for the
ATE / CATE estimator family from the econometrics literature,
\pkg{dowhy} (Microsoft) for end-to-end causal-inference
pipelines, and \pkg{CausalPy} (PyMC) for Bayesian
quasi-experimental designs. The package does not duplicate these
established estimators where they already exist in the
\proglang{Python} stack; instead, it wraps them under a single
function-naming convention and pairs them with the MRM modules
that target Canadian open carceral and police-oversight data.

The MRM modules implement the unified framework of
\citet{Ruhela2026MRM} for analysis across five Canadian sources:
the Ontario Tracking Information System (\OTIS, A01-RCDD), the
federal Structured Intervention Unit Implementation Advisory
Panel, the Ontario Special Investigations Unit (\SIU), the
Toronto Police Service (\TPS) open data, and the federal
Corrections and Conditional Release Statistical Overview. The
framework's theoretical foundations are presented in
\citet{Ruhela2026MRM}; this paper concerns its \proglang{Python}
software realisation.

The companion paper \citep{Ruhela2026MorieR} describes the
\proglang{R} implementation, which is functionally parallel to the
\proglang{Python} implementation under a different idiomatic API.

The rest of the paper is organised as follows.
Section~\ref{sec:getting-started} demonstrates installation and a
motivating example. Section~\ref{sec:taxonomy} surveys the module
taxonomy. Section~\ref{sec:mrm} describes the MRM modules.
Section~\ref{sec:api} contrasts the functional and object-oriented
APIs and documents the \pkg{RichResult} return type.
Section~\ref{sec:wrapper} documents the wrapper architecture
around \pkg{DoubleML}. Section~\ref{sec:examples} works three
reproducible examples.  Section~\ref{sec:mrm-empirical-callables}
catalogues the twelve empirical-workflow callables added in v0.1.15.
Section~\ref{sec:textbook-callables} catalogues the 275
textbook-derived callables shipped in \code{morie.fn} across
fourteen suites in v0.2.1 and v0.3.0.
Section~\ref{sec:simulation} reports a simulation study.
Section~\ref{sec:conclusion} concludes.

\section{Getting started}
\label{sec:getting-started}

\subsection{Installation}

The released \pkg{morie} package is available on PyPI:

\begin{Code}
$ pip install morie
\end{Code}

\noindent Every wheel and source distribution on PyPI ships with a
PEP 740 sigstore attestation; the provenance can be verified via
the PyPI integrity endpoint:

\begin{Code}
$ pip install pip-audit
$ pip-audit --strict morie
\end{Code}

\noindent The development version is installable from GitHub:

\begin{Code}
$ pip install "git+https://github.com/hadesllm/morie"
\end{Code}

\noindent For users who lack \pkg{pip} or want a self-contained
managed Python environment, a one-line installer is published on the
documentation site.  It bootstraps \pkg{uv}, materialises a Python
3.12 virtual environment, and installs the wheel inside it:

\begin{Code}
$ curl -fsSL https://hadesllm.github.io/morie/install.sh | bash
\end{Code}

\noindent A Homebrew tap is available for \proglang{macOS} and
Linuxbrew users:

\begin{Code}
$ brew tap hadesllm/morie
$ brew install morie
\end{Code}

\noindent The tap source is at
\href{https://github.com/hadesllm/homebrew-morie}{hadesllm/homebrew-morie}
and bundles a self-contained \texttt{python@3.12} virtualenv with the
SciPy stack.

\noindent The package requires \proglang{Python} 3.10 or later
and declares \pkg{pandas}, \pkg{numpy}, \pkg{scipy},
\pkg{scikit-learn}, \pkg{statsmodels}, and \pkg{DoubleML} as
hard dependencies. The interactive TUI components, the carbon
emissions tracker, and the imbalanced-learn integration are
declared under optional extras to keep the base installation
lightweight.

\subsection{A motivating example}

The motivating example uses the open \OTIS{} A01-RCDD record to
estimate the average treatment effect of an indicator of high
alert complexity on placement-volatility outcomes in the Ontario
restrictive-confinement system.

\begin{Code}
>>> import morie
>>> data = morie.otis_load_a01_rcdd(year=2025)
>>> data.head()
\end{Code}

\noindent The treatment variable is constructed by thresholding
the alert-count covariate at the population median. The outcome is
a count of placement transitions within the fiscal year.

\begin{Code}
>>> data["D"] = (data["alert_count"] >= data["alert_count"].median()).astype(int)
>>> result = morie.causal.estimate_irm(
...     y="placement_transitions",
...     d="D",
...     x=["age_band", "gender", "regional_cluster"],
...     data=data,
... )
>>> print(result)
\end{Code}

\noindent The \code{estimate\_irm()} wrapper composes a
\pkg{scikit-learn} regression estimator
(\code{LinearRegression()}) and a classification estimator
(\code{LogisticRegression()}) and passes them to
\code{DoubleML.DoubleMLIRM}. The returned object is a
\pkg{RichResult} (Section~\ref{sec:api}) that prints a
multi-section interpretation alongside the point estimate,
standard error, and confidence interval.

\section{Module taxonomy}
\label{sec:taxonomy}

The \pkg{morie} \proglang{Python} interface organises
functionality into eight thematic submodules.

\subsection{morie.causal}

The causal submodule provides the average-treatment-effect
estimator family:
\code{morie.causal.estimate\_att}, \code{estimate\_atc},
\code{estimate\_aipw}, \code{estimate\_cate},
\code{estimate\_gate}, \code{estimate\_late},
\code{estimate\_irm}, and \code{estimate\_plr}. The first four
are implemented in pure \proglang{Python} using the identifying
assumptions described in \citet{Ruhela2026MRM}; the latter four
wrap \pkg{DoubleML} as detailed in Section~\ref{sec:wrapper}.

\subsection{morie.weights}

The weights submodule provides
\code{morie.weights.calculate\_ipw\_weights} (the propensity-score
weighting routine with stabilisation and trimming),
\code{morie.weights.calculate\_ebac} (the Widmark
blood-alcohol-content estimator used by Health Canada survey
instruments), and helpers for design-based variance estimation
via \pkg{statsmodels}.

\subsection{morie.signal}

The signal submodule provides Butterworth filter primitives
(\code{buttlp}, \code{butthp}, \code{buttbp}, \code{buttbs}),
Savitzky--Golay smoothing (\code{sgolay\_smooth}), Hurst-exponent
estimation by R/S analysis (\code{hurst\_r}), Higuchi fractal
dimension (\code{hfd}), and a phonocardiogram band-pass filter
(\code{pcg\_filter}).

\subsection{morie.spatial}

The spatial submodule provides Moran's $I$ and Geary's $C$
spatial autocorrelation estimators, kernel density estimators,
and the Ripley-$K$ wrapper used by the MRM spatio-temporal
modules.

\subsection[morie.tps stochastic]{morie.tps\_stochastic}

This submodule provides Hawkes process maximum-likelihood
estimation (\code{morie.tps\_stochastic.fit\_hawkes}; exponential
Weibull and Lomax kernels under both constant and sinusoidal baselines), reaction--diffusion intensity-field
estimation, L\'evy-flight displacement modelling, and
Lotka--Volterra police--crime dynamics
\citep{Mohler2011, Ruhela2026Hawkes}.

\subsection{morie.psychometrics}

The psychometrics submodule provides classical-test-theory
estimators (\code{cronbach\_alpha}, \code{mcdonald\_omega},
\code{kuder\_richardson}) and item-response-theory estimators
(1PL, 2PL, 3PL Rasch family).

\subsection{morie.crypto}

The cryptography submodule provides reference implementations of
ML-KEM, ML-DSA, NTRU, and hybrid post-quantum primitives. These
are intended for educational use and are not audited for
production deployment; the package's \texttt{SECURITY.md}
documents this status explicitly.

\subsection{morie.investigation}

The investigation submodule provides
\code{compare\_nested\_logistic\_models} (likelihood-ratio test
across nested specifications),
\code{run\_treatment\_effects\_analysis} (end-to-end
ATE/ATT/CATE/GATE estimation pipeline),
\code{run\_weighted\_logistic\_analysis} (survey-aware logistic
regression), and \code{inspect\_output} (audit-trail inspection
of result objects).

\section{The MRM modules}
\label{sec:mrm}

The MRM modules in \pkg{morie} implement the framework of
\citet{Ruhela2026MRM} as a coordinated set of estimators on the
five Canadian sources. Each MRM module is exposed by an
\code{mrm\_*} prefixed function at the package top level:
\code{morie.mrm\_otis\_load}, \code{morie.mrm\_siu\_scrape},
\code{morie.mrm\_tps\_load}, \code{morie.mrm\_classify\_mandela},
\code{morie.mrm\_chi2\_doob}, and \code{morie.mrm\_hawkes\_tps}.
The functions handle source-specific schema reconciliation
through the canonical fiscal-quarter index of
\citet{Ruhela2026MRM} and emit \pkg{pandas} \code{DataFrame}
objects that the estimator-family functions of
Section~\ref{sec:taxonomy} consume directly.

\section{Functional and object-oriented APIs}
\label{sec:api}

\pkg{morie} exposes both a functional API (the
\code{morie.causal.estimate\_*} family of free functions) and an
object-oriented API (a single \code{Pipeline} class that
internally orchestrates estimator selection, nuisance estimation,
and result reporting). The functional API is the canonical entry
point for one-shot analyses; the object-oriented API is provided
for the composable analysis pipelines that benefit from inspection
of intermediate state.

\subsection{The RichResult return type}

Every estimator in \pkg{morie} returns a \code{RichResult}, a
dataclass-derived object that inherits from \code{dict} and
populates itself with the canonical result fields
(\code{estimate}, \code{se}, \code{ci\_lower}, \code{ci\_upper},
\code{p\_value}, \code{n}, \code{n\_treated}, \code{interpretation}).
The dataclass inheritance pattern provides two simultaneous use
modes: \code{isinstance(result, dict)} returns \code{True} (so
results can be passed to any \code{dict}-consuming downstream
code), and \code{result.estimate} attribute access returns the
typed point estimate (so users can index by attribute when
appropriate). The multi-section pretty-printer renders the result
as a panel of labelled fields when called interactively; when
used programmatically, the result behaves as an ordinary
dictionary.

\subsection{describe()}

Every callable in \pkg{morie} that returns a \code{RichResult}
has a paired \code{describe\_<callable>.md} pedagogical narrative
shipped as a package data file. The \code{morie.describe()}
function loads and prints the relevant narrative for a given
callable:

\begin{Code}
>>> morie.describe(morie.causal.estimate_irm)
\end{Code}

\noindent The describe-output convention is documented in the
package's design paper at the repository.

\section{Wrapper architecture for DoubleML}
\label{sec:wrapper}

\pkg{morie} layers a thin function-style interface over
\pkg{DoubleML}'s class-based interface. Where \pkg{DoubleML}
exposes \code{DoubleMLIRM(data, ml\_g, ml\_m, ...)} with a data
backend and two \pkg{scikit-learn} estimators as separate
arguments, \code{morie.causal.estimate\_irm} accepts a single
\pkg{pandas} \code{DataFrame}, column names for the treatment,
outcome, and covariates, and an optional learner specification
with sensible defaults.

The default learners for \code{estimate\_irm} are
\code{LinearRegression()} (linear regression for the outcome
nuisance) and \code{LogisticRegression()} (logistic regression
for the propensity nuisance), chosen for the absence of
additional dependencies. Users who require flexible non-parametric
nuisance estimation pass alternative \pkg{scikit-learn} estimators
directly:

\begin{Code}
>>> from sklearn.ensemble import RandomForestRegressor, RandomForestClassifier
>>> result = morie.causal.estimate_irm(
...     y="y", d="D", x=["X1", "X2", "X3"],
...     data=df,
...     ml_g=RandomForestRegressor(n_estimators=500),
...     ml_m=RandomForestClassifier(n_estimators=500),
...     n_folds=5, n_rep=2,
... )
\end{Code}

\noindent The wrapper performs three additional services beyond
the \pkg{DoubleML} class call. First, it harmonises the result
object to a \code{RichResult}. Second, it audits learner-result
quality through cross-validation $R^{2}$ and Brier-score
diagnostics. Third, it surfaces both the ATE and the ATT in a
single call rather than requiring a separate \code{score = "ATTE"}
parameterisation.

The wrapper preserves the full \pkg{DoubleML} class instance under
\code{result["doubleml"]} for users who require direct access to
sample splits, bootstrap statistics, or the underlying psi-score
arrays.

\section{Reproducible examples}
\label{sec:examples}

Three reproducible examples illustrate the package's use across
the MRM source set.

\subsection[Example 1: alert complexity and placement volatility (OTIS)]{Example 1: alert complexity and placement volatility (\OTIS{})}

The first example revisits the motivating example of
Section~\ref{sec:getting-started} under three estimators
(IPW, AIPW, IRM-DML) and reports the consistency of point
estimates across estimators as evidence of identification
robustness.

\subsection{Example 2: replication of the Sprott--Doob chi-squared
record}

The second example replicates the published $\chi^{2}$ statistics
from \citet{SprottDoob2023} on the federal $\SIU$ aggregate
record. The MRM \code{mrm\_chi2\_doob} function consumes the
published contingency tables and reproduces every reported
$\chi^{2}$ statistic to within $0.01$.

\subsection{Example 3: Hawkes self-exciting process on TPS major
crime}

The third example fits the Hawkes self-exciting specification of
\citet{Ruhela2026Hawkes} to the Toronto Police Service open-data
major-crime event series.  The Markovian (exponential kernel,
constant baseline) and a fat-tailed non-Markovian (Lomax /
power-law kernel, sinusoidal baseline) specifications are fit
by maximum likelihood under the Numba-JIT'd $O(n)$ recursive
likelihood for the exponential kernel and the parallel $O(n^2)$
JIT path for non-Markovian kernels (new in \pkg{morie} v0.2.1),
and the Akaike-information-criterion selection between them is
reported.  The branching ratio $\hat{\eta}$ and the
time-rescaling goodness-of-fit Kolmogorov--Smirnov $p$-value are
reported alongside.  On the canonical TPS Assault 2019
contiguous one-year window ($n=21{,}160$ events), the Markovian
fit attains $\hat\eta=0.567$, $\hat\beta=3.854$/day, and
KS\,$p=0.494$; the Exp/sin model improves AIC by $27.7$
(decisive on the standard $\Delta\mathrm{AIC} > 10$ rule) to
$\hat\eta=0.510$, KS\,$p=0.332$; and Lomax/sin lands at
$\hat\eta=0.513$, KS\,$p=0.298$ --- $3.9$ AIC \emph{worse} than
Exp/sin (one extra parameter, no likelihood gain).  The
substantive finding for TPS Assault 2019 is that
non-stationarity (sinusoidal baseline) dominates
non-Markovianity (kernel shape); a richer kernel does not pay
its parameter cost on this data.

\section{Simulation study}
\label{sec:simulation}

A simulation study under the partially-linear data-generating
process of \citet{Chernozhukov2018DML} compares the
\code{estimate\_irm} wrapper's coverage and bias under three
nuisance-learner choices: linear (\code{LinearRegression},
default), lasso (\code{LassoCV}), and random-forest
(\code{RandomForestRegressor}). The simulation uses $n = 500$
observations and $p = 20$ covariates. The true treatment effect
is $\theta_0 = 0.5$. Each setting is replicated $500$ times.

Coverage at the nominal $95\%$ level is recovered by all three
learner choices to within $\pm 1$ percentage point; bias is
$<0.01$ for the linear and lasso learners and $<0.03$ for the
random-forest learner under default hyperparameters.

\section[Empirical workflow callables added in v0.1.15]{Empirical workflow callables added in \texttt{v0.1.15}}\label{sec:mrm-empirical-callables}

Version 0.1.15 ships twelve empirical-workflow callables on
the \proglang{Python} side with full R parity.  The canonical
enumeration of these callables --- with method, data source,
and verification target per callable --- lives in the
companion \proglang{R} paper \citep[Section 8]{Ruhela2026MorieR}.
Briefly, the \proglang{Python} side exposes the same
twelve names under \code{morie.mrm\_otis\_*},
\code{morie.mrm\_tps\_*}, \code{morie.mrm\_siu\_*}, plus
the longitudinal-panel simulator and dataset fetchers.  Bit
identical numeric agreement between the \proglang{R} and
\proglang{Python} sides is verified by the package smoke
suite (e.g. mortification-co-occurrence Cramér's $V=0.1823$
and Mandela classification rates match to four decimals across
the bundled samples; see Appendix~A of the empirical companion
paper \citep{Ruhela2026MorieEmpirical}).  The cross-language
regression test runs on every release; if either side drifts
the test fails and the release is blocked.

\section[Textbook-derived callables in morie.fn]{Textbook-derived callables in \texttt{morie.fn}}\label{sec:textbook-callables}

Beyond the causal-inference core and the empirical-workflow
callables of Section~\ref{sec:mrm-empirical-callables},
\pkg{morie} ships a curated catalogue of 275 textbook-derived
callables, organised into fourteen thematic suites under the
\code{morie.fn} namespace.  Each suite implements a coherent
slice of a canonical reference text in statistics, machine
learning, or signal processing; together they form a
pedagogical-and-applied corpus that complements the wrappers
documented above.  The design philosophy is uniform: where a
battle-tested library exists on either side of the
\proglang{Python}/\proglang{R} divide, the suite calls into it
(\code{scipy}, \code{statsmodels}, \code{scikit-learn},
\code{pywavelets}, \code{arch} on the \proglang{Python} side;
\code{spdep}, \code{gstat}, \code{survival}, \code{forecast},
\code{rugarch}, \code{signal}, \code{kernlab} on the
\proglang{R} side); where the operation is light enough that a
library would introduce more weight than logic, the suite uses
\code{numpy} primitives (or base \proglang{R} matrix
primitives on the sibling side) directly.  All 275 callables
were verified for \proglang{Python} + \proglang{R} parity at the
v0.3.0 release (see the package \code{NEWS} file).

The fourteen suites are enumerated in
Table~\ref{tab:textbook-suites}, with the short submodule
mnemonic that fixes the namespace, the number of callables
shipped, the canonical reference textbook, and one representative
fully-qualified \proglang{Python} call as illustration.

\begin{table}[t]
\centering
\scriptsize
\begin{tabular}{lr>{\raggedright\arraybackslash}p{4.5cm}>{\raggedright\arraybackslash}p{5.8cm}}
\toprule
Suite & $n$ & Reference & Representative call \\
\midrule
\code{schabenberger} & 20 & \citet{SchabenbergerGotway2005} & \code{morie.fn.sptau.spatial\_autocorrelation()} \\
\code{kosorok} & 20 & \citet{Kosorok2008} & \code{morie.fn.ksr01.kosorok\_empirical\_process()} \\
\code{ml\_foundations} & 20 & \citet{MohriRostamizadehTalwalkar2018,HastieTibshiraniFriedman2009} & \code{morie.fn.linrg.linear\_regression\_ols()} \\
\code{ghosal} & 20 & \citet{GhosalVanDerVaart2017} & \code{morie.fn.ghdir.ghosal\_dirichlet\_posterior()} \\
\code{horowitz} & 20 & \citet{Horowitz2009} & \code{morie.fn.hrzk2.horowitz\_kernel\_regression()} \\
\code{montesinos} & 20 & \citet{MontesinosLopez2022} & \code{morie.fn.gblpf.gblup\_full()} \\
\code{armstrong} & 20 & \citet{Armstrong2014SpatialVoting} & \code{morie.fn.idlpt.ideal\_point\_recovery()} \\
\code{deep\_learning} & 20 & \citet{GoodfellowBengioCourville2016} & \code{morie.fn.fwpas.forward\_pass\_dense()} \\
\code{missing\_stats} & 20 & QMC / copulas / EVT & \code{morie.fn.btsrp.bootstrap\_ci()} \\
\code{llm\_arch} & 20 & \citet{Vaswani2017Attention,Dao2022FlashAttention,Su2021RoPE,ZhangSennrich2019RMSNorm} & \code{morie.fn.tknbp.bpe\_tokenizer()} \\
\code{fauzi} & 15 & \citet{Fauzi2018Kernel} & \code{morie.fn.fzkdf.fauzi\_kdfe\_properties()} \\
\code{gibbons\_chakraborti} & 20 & \citet{GibbonsChakraborti2014} & \code{morie.fn.tolim.tolerance\_limits()} \\
\code{rangayyan} & 20 & \citet{Rangayyan2015} & \code{morie.fn.rgfir.rangayyan\_fir\_filter()} \\
\code{time\_series\_advanced} & 20 & \citet{HyndmanAthanasopoulos2018,Tsay2010FinancialTimeSeries} & \code{morie.fn.garch.garch\_fit()} \\
\bottomrule
\end{tabular}
\caption{The fourteen textbook-derived suites in \code{morie.fn}.
Total: 275 callables, all with \proglang{Python}+\proglang{R}
parity verified at v0.3.0.  The \code{missing\_stats} label
is a v0.2.1 artefact: the spec covers resampling,
quasi-Monte-Carlo, copula, and extreme-value-theory material,
not missing-data methods proper; a true missing-data spec is
deferred to v0.3.x.}
\label{tab:textbook-suites}
\end{table}

\subsection{Naming convention}\label{sec:textbook-naming}

Each suite occupies a four-to-seven-character mnemonic submodule
under \code{morie.fn}.  The \proglang{Python} fully-qualified
form is \code{morie.fn.<short>.<full\_name>(\ldots)}, where the
mnemonic matches the spec \code{name} field of the originating
formula-index JSON and the function name matches the
\code{full\_name} field, so a user can locate the originating
specification and the cross-language sibling without guessing.  The \proglang{R} sibling additionally aliases the
mnemonic so that both \code{sptau()} and
\code{spatial\_autocorrelation()} resolve in a single namespace,
as documented in \citet{Ruhela2026MorieR}.

\subsection{Caveats and v0.3.x deferrals}\label{sec:textbook-caveats}

Numerical agreement between the \proglang{Python} and
\proglang{R} sides is bit-identical for deterministic callables:
the closed-form OLS and PCA in \code{ml\_foundations}, Moran's
$I$ and ordinary kriging in \code{schabenberger}, the analytic
attention and softmax operators in \code{llm\_arch}, and similar.  For stochastic callables
(bootstrap and jackknife in \code{missing\_stats}, random-forest
and gradient-boosting ensembles, MCMC posteriors in
\code{ghosal}, the Bayesian-inference paths in
\code{time\_series\_advanced}), agreement is within Monte-Carlo
tolerance under matched random-number seeds.  The
\proglang{Python} smoke runs the \code{\_\_main\_\_} canonical-test
block in every shipped module across all 275 callables; the
\proglang{R} smoke confirms that all 319 \proglang{R} files in the
source tree source cleanly under base \proglang{R} with only the
declared \code{Suggests:} optional packages installed.  Three
honest deferrals carry over from the v0.2.1 tracker: an opt-in
deterministic mode for the stochastic callables, a true
missing-data suite to replace the misnamed \code{missing\_stats}
spec, and a tighter cross-language MLE-tolerance pin for the
GARCH/EGARCH and Bayesian fits where the two language
ecosystems' optimiser defaults differ.

\section{Conclusion}
\label{sec:conclusion}

The \proglang{Python} package \pkg{morie} provides a curated
multi-domain toolkit for observational inference, with explicit
support for the MRM framework as a sociolegal application target.
The wrapper architecture around \pkg{DoubleML} offers a thin
function-style API; the eight thematic submodules provide a
unified namespace for the broader observational-inference,
signal-processing, spatial, statistical-physics, and psychometric
routines that complement the causal core. The MRM modules
implement the framework of \citet{Ruhela2026MRM} on the five
Canadian sources through a common ingestion pipeline.

The package is released on PyPI with PEP 740 sigstore
attestations on every wheel and source distribution, and the
source repository is at \url{https://github.com/hadesllm/morie}.
Contributions follow the process documented in the
\texttt{CONTRIBUTING.md} of the source repository.

\section*{Acknowledgements}

The author thanks Glenn McNamara for distribution-theory
mentorship that grounds the per-individual estimator ensemble;
Professor Angela Zorro Medina for the methodological lineage of
the framework's identification strategy; and the DoubleML team
(Philipp Bach, Malte S.\ Kurz, Victor Chernozhukov, Martin
Spindler, Sven Klaassen) for the BSD-3-Clause
double-machine-learning library on which the package's IRM and
PLR wrappers are layered.

\bibliography{refs}

\end{document}
